\section{Pronombres personales (singular y plural)}
\textbf{Objetivo:} Identificar y usar los pronombres sujeto en todas las personas.

\textbf{Lista rápida}
\begin{itemize}[leftmargin=*]
  \item Singular: \textit{ik} (yo), \textit{je/jij} (tú), \textit{u} (usted), \textit{hij} (él), \textit{zij/ze} (ella), \textit{het} (ello/neutral).
  \item Plural: \textit{we/wij} (nosotros), \textit{jullie} (vosotros/ustedes informal), \textit{zij/ze} (ellos/ellas).
\end{itemize}

\textbf{Notas de uso}
\begin{itemize}[leftmargin=*]
  \item \textit{jij} enfatiza, \textit{je} es neutro; \textit{u} es formal.
  \item \textit{het} se usa con objetos/neutros; en la práctica oral puede omitirse si ya se nombró.
\end{itemize}

\textbf{Práctica}
\begin{itemize}[leftmargin=*]
  \item Sustituye el sujeto en frases dadas para recorrer todas las personas.
  \item Escucha/lee diálogos marcando cada pronombre.
\end{itemize}
